\section{Introduction} \label{sec_intro}
The t-distributed stochastic neighbor embedding(t-SNE) is a well-known method of visualizing high-dimensional data with low-dimensional vectors (so they can be easily visualized). Let $d,s \geq 2$ with $d>s$, and $n \geq 1$. The input of t-SNE consist of vectors $X_1, \cdots, X_n \in \mathbb{R}^d$. The goal is the find vectors $Y_1, \cdots Y_n \in \mathbb{R}^s$ that serve as a good representation of the original vectors. The process of t-SNE starts with defining the following probability masses:
\begin{align} \label{original_prob_mass}
    p_{j|i} = \frac{\exp(-\norm{X_i-X_j}^2/2\sigma_i^2)}{\sum_{k \neq i} \exp(-\norm{X_i-X_k}^2/2\sigma_i^2)},
\end{align}
where $\sigma_i$ are given constants. For each pair of data points $X_i$ and $X_j$, we assign a weight to that pair, namely the Gaussian function $\exp(-\norm{X_i-X_j}^2/2\sigma_i^2)$. We call such function the \textit{input kernel}. In this paper, we are looking at generalized kernels. Let $\ki(x,y,\sigma): \mathbb{R}^d  \times \mathbb{R}^d \to \mathbb{R}^+$ be such that $\ki(x,y,\sigma)=\ki(y,x,\sigma)$. We replace the probability mass in (\ref{original_prob_mass}) with the following
\begin{align}
    p_{j|i} = \frac{\ki(X_i,X_j,\sigma_i)}{\sum_{k \neq i} \ki(X_i,X_k,\sigma_i)}.
\end{align}
$\sigma_i$ satisfies that for each $i = 1, \dots, n$, that
\begin{align}
    -\sum_{j = 1, j \neq i}^n p_{j|i} \log p_{j|i} = \log \textbf{Prep} =: \log (n \rho).
\end{align}
$\textbf{Prep}$ is called the perplexity of the model and $\rho \in (0,1)$ is a given constant. The purpose of $\sigma_i$ is to make sure that for each $i$, the probability distribution induced by $p_{j|i}$ has the same entropy across all $i$. Clearly, $p_{j|i}$ is not symmetric. We then define $p_{ii} = 0$ and for $1 \leq i \neq j \leq n$,
\begin{align}
    p_{ij} = \frac{p_{j|i} + p_{i|j}}{2n} = \frac{1}{2n}\left(\frac{\ki(X_i,X_j,\sigma_i)}{\sum_{k \neq i} \ki(X_i,X_k,\sigma_i)}+\frac{\ki(X_i,X_j,\sigma_j)}{\sum_{k \neq j} \ki(X_j,X_k,\sigma_j)}\right).
\end{align}
Obviously,
\begin{align}
    \sum_{\substack{j = 1 \\ j \neq i}}^n p_{j|i} = 1, \quad \sum_{1 \leq i \neq j \leq n} p_{ij} = 1.
\end{align}
As for the output vectors $Y_1, \cdots Y_n \in \mathbb{R}^s$, set for $1 \leq i \neq j \leq n$,
\begin{align}
    q_{ij} = \frac{\ko(Yi,Yj)}{\sum_{k \neq l}\ko(Y_k,Y_l)}.
\end{align}
The weight we assign to each pair of output data points $Y_i$ and $Y_j$, a.k.a. a function $\ko(Y_i,Y_j):\mathbb{R}^s \times \mathbb{R}^s \to \mathbb{R}^+$, is called the output kernel. In the original algorithm, $\ko(Y_i,Y_j) = (1+\norm{Y_i-Y_j}^2)^{-1}$. Denote $X = (X_1,\cdots,X_n)$ and $Y = (Y_1,\cdots,Y_n)$ the input and output data set. The goal of t-SNE is to find $Y$ that minimizes the relative entropy
\begin{align}
    L_{n,\rho}(X,Y) = \sum_{1 \leq i \neq j \leq n} p_{ij} \log \left(\frac{p_{ij}}{q_{ij}}\right).
\end{align}
More precisely, the output of t-SNE $Y^*$ is defined as
\begin{align}
    Y^* = \argmin_{Y \in (\mathbb{R}^s)^n}L_n(X,Y).
\end{align}
In the expression of $L_n(X,Y)$, all $p_{ij}$ are given and $q_{ij}$ is a function of $Y_i$ and $Y_j$. Therefore, $Y^*$ can be found simply by using gradient descent. This is no different from the original version of t-SNE. However, it is the mathematical properties and implications of t-SNE that we are interested in. Towards this end, let $\mathcal{M}(\mathbb{R}^d)$ be the space of Borel probability measures on $\mathbb{R}^d$ and $\mu_X \in \mathcal{M}(\mathbb{R}^d)$ be a probability measure. Suppose all input $X = (X_1,\cdots,X_n)$ are i.i.d., independent random variables with probability measure $\mu_X$. Meanwhile, let $Y^* = (Y_1,\cdots,Y_n)$ be the output of the t-SNE algorithm. Consider the following empirical measure:
\begin{align}
    \mu_n := \frac{1}{n} \sum_{i=1}^n \delta(X_i,Y_i) \in \mathcal{M}(\mathbb{R}^d \times \mathbb{R}^s).
\end{align}
We are interested in its limiting measure as $n \to \infty$. To finish the setup, there are a few more remaining quantity we would like to introduce. Let $\mu \in \mathcal{M}(\mathbb{R}^d \times \mathbb{R}^s)$. For a given kernel $\ki(x,y)$, define $F_{\rho,\mu}(x,\sigma) : \mathbb{R}^d \times \mathbb{R}^+ \to \mathbb{R}$ as
\begin{align}
    F_{\rho,\mu}(x,\sigma) = \int \frac{\ki(x,x',\sigma)}{\int \ki(x,x',\sigma)d\mu(x')} \log \left(\frac{\ki(x,x',\sigma)}{\int \ki(x,x',\sigma)d\mu(x')}\right) d\mu(x') + \log \rho.
\end{align}
Note that the last $s$ variables in the integrands stay constant. It is not hard to see that $F_{\rho,\mu}(x,\sigma)$ is the continuous version of the entropy. Next, we extend the definition of $\sigma_i$ by defining $\sigma^*_{\rho,\mu} : \mathbb{R}^d \to \mathbb{R}$ the unique solution of the equation
\begin{align}
    F_{\rho,\mu}(x,\sigma) = 0.
\end{align}
The introduction of $\sigma^*$ aligns with the procedure of finding $\sigma_i$ that controls the entropy to be the same for all $i$. Of course, the existence and uniqueness of $\sigma^*$ is not trivial and will be discussed in section \ref{tsne:sec:existence_uniqueness}. For the symmetric probability mass, let $\psi: \mathbb{R}^d \to \mathbb{R}$ be a given function and $p_{\psi}: \mathbb{R}^d \times \mathbb{R}^d \to \mathbb{R}$ be defined by
\begin{align}
    p_{\psi}(x,x') = \frac{1}{2}\left( \frac{\ki(x,x',\psi(x))}{\int \ki(x,x',\psi(x))d\mu(x')} + \frac{\ki(x,x',\psi(x'))}{\int \ki(x,x',\psi(x'))d\mu(x)}\right).
\end{align}
For the extension on the output side, let $q: \mathbb{R}^s \times \mathbb{R}^s \to \mathbb{R}$ be given by
\begin{align}
    q(y,y') = \frac{\ko(y,y')}{\iint \ko(y,y')d\mu(y) d\mu(y')}.
\end{align}
Based on the construction of $p$ and $q$, define
\begin{align}
    I_{\rho}(\mu) = \iint p_{\sigma^*_{\rho,\mu}}(x,x') \log\left(\frac{p_{\sigma^*_{\rho,\mu}}(x,x')}{q(y,y')}\right) d\mu(x,y) d\mu(x',y').
\end{align}
Recall $\mu_X \in \mathcal{M}(\mathbb{R}^d)$ is an arbitrary given probability measure. We set $\mathcal{P}_X$ to be the set of probability measures on $\mathbb{R}^d \times \mathbb{R}^s$ whose marginal coincides with $\mu_X$. Formally speaking,
\begin{align}
    \mathcal{P}_X := \{\mu \in \mathcal{M}(\mathbb{R}^d \times \mathbb{R}^s): \text{the $\mathbb{R}^d$ marginal of $\mu$ is $\mu_X$}\}.
\end{align}
Lastly, define
\begin{align}
    \tilde{\mathcal{P}}_X := \left\{ \mu \in \mathcal{P}_X: \int (p_{\sigma^*_{\rho,\mu}}(x,x')-q(y,y'))\frac{\frac{\partial}{\partial y}\ko{(y,y')}}{\ko{(y,y')}} d\mu(x',y') = 0,\mu(x,y)\text{-a.e.}\right\}.
\end{align}
% The goal of this project is to find the relationship between the following two quantities:
% \begin{align}
%     \lim_{n \to \infty} \inf_Y L(X,Y), \text{ and} \inf_{\mu \in \tilde{\mathcal{P}}_X} I(\mu).
% \end{align}
\section{Main Results}
The core results of this problem are the conditions we assume the input and output kernels to have such that the output of the generalized $t$-SNE algorithm converges to a limiting equilibrium measure. Towards this end, below are the constraints we impose on the input and output kernels.
\begin{condition}[Input Kernel Conditions] \label{input_kernel_conditions}
    Let $w: \mathbb{R}_{\geq 0} \to \mathbb{R}_{\geq 0}$ be some fixed function. The input kernel has the form of $\ki(x_i,x_j,\sigma_i) = \exp(-w(\sigma_i\norm{x_j-x_i}^{\theta}))$. Here, $\sigma_i$ is used to control entropy (thus perplexity) and $\theta>0$ is a fixed constant. (We can replace $\sigma_i$ with $\psi: \mathbb{R}^d \to \mathbb{R}^+$ in the continuous version of the kernel.) Define
    \begin{align}
      \mathcal{K}_w = \left\{\tilde{C} \in \mathbb{R}_{\geq 0} \bigg\rvert \lim_{n \to \infty} n^{-1/2} (\log n)^{\theta} \exp \left(4 w\left(\left(\tilde{C} \sqrt{\log n} \right)^{\theta} \right) \right) = 0 \right\}.
    \end{align}
    Then $w$ satisfies that
    $w$ satisfies:
  
    1. $\forall t \geq 0$, $w'(t)>0$ and $\exists M > 0$ s.t. $w'(t) \leq M$. $\lim_{t \to \infty} w(t) = \infty$;
  
    2. $w'(t)+tw''(t) \geq 0$ for all $t \geq 0$;
  
    3. The following integral converges:
    \begin{align}
      \int_0^{\infty} t^{d-1} w(t^{\theta}) \exp(-w(t^{\theta}))dt < \infty;
    \end{align}
  
    4. $\sup \mathcal{K}_w > 0$.
\end{condition}
\begin{condition}[Output Kernel Conditions] \label{output_kernel_conditions}
    The output kernel $\ko(y,y') = k (\norm{y-y'})$ depends only on the distance between $y$ and $y'$ and that:
  
    1. $k: \mathbb{R}^+ \to \mathbb{R}^+$ is decreasing and bounded;
  
    2. $g$ satisfies
    \begin{align}
      \iint_{(0,\infty)^2} g(y,y') dy dy' < \infty;
    \end{align}
  
    3. $k$ has bounded derivative;
  
    4. $k'(0) = 0$.
\end{condition}
We are ready to state the main results
\begin{theorem} \label{tsne:main_theorem_1}
    Assume $\mu_X$ has a $C^1$ density $f(x)$ and compact support. Let $X_i \in \mathbb{R}^d$ be i.i.d random variables with probability measure $\mu_X$. Then for any $\rho \in (0,1)$,
    \begin{align}
        \lim_{n \to \infty} \inf_Y L_{n,\rho}(X,Y) = \inf_{\mu \in \tilde{\mathcal{P}}_X} I_{\rho}(\mu).
    \end{align}
    Moreover, there exist a measure $\mu^*$ and a sub-sequence $\{n_k\}$ such that $\mu_{n_k} \to \mu^*$ weakly and $I_{\rho}(\mu^*) = \inf_{\mu \in \tilde{\mathcal{P}}_X} I_{\rho}(\mu)$.
\end{theorem}
\begin{theorem}
    Under the same assumptions of theorem \ref{tsne:main_theorem_1}, $\mu^*$ has compact support.
\end{theorem}
\section{Intuitive Insights and Simulations}
\gy{TODO: Come back and finish this section when the rest of the paper is complete.}
\begin{figure}
    \centering
    \subfloat[\centering Original $t$-SNE]{{\includegraphics[width=6.5cm]{default_kernel.png} }}%
    \qquad
    \subfloat[\centering Polynomial Input Kernel]{{\includegraphics[width=6.5cm]{poly_input_kernel.png} }} \\
    \subfloat[\centering Mixed Input Kernel]{{\includegraphics[width=6.5cm]{poly+exp_input.png} }}%
    \qquad
    \subfloat[\centering Custom Input\&Output Kernels]{{\includegraphics[width=6.5cm]{custom_in+out.png} }}
    \caption{$t$-SNE running with different combinations of input and output kernels on the MNIST data set. All three custom kernel combinations are able to perform on par with the original algorithm.}%
    \label{fig:tsne_plots}
  \end{figure}
\clearpage
\section{Existence and Uniqueness of \texorpdfstring{$\sigma^*$}{Lg}}\label{tsne:sec:existence_uniqueness}
\section{Uniform Convergence of \texorpdfstring{$\sigma^*$}{Lg}}\label{tsne:sec:uniform_convergence}
\section{Proofs of the Main Results}\label{tsne:sec:proofs}
\section{Extra Technical Proofs}\label{tsne:sec:extra_proofs}